\section{Introduction}
\label{sec:intro}

Conditional computation is usually motivated by the same accounting argument:
most tokens do not need the whole model, so spend less on them. The argument is
made in two places, and the two answer different questions. Sparse
mixture-of-experts (MoE) layers \citep{jacobs1991moe,shazeer2017moe,fedus2022switch}
decide \emph{which} computation module processes a token, dispatching it to one
of several independent sub-networks. Recursive weight-tied architectures with
learned halting \citep{graves2016act,dehghani2019ut,banino2021pondernet} decide
\emph{how much} computation a token receives, reusing one block for a variable
number of steps.

Because the two decisions are made over different objects---which parameters
versus how many applications of them---they are orthogonal in principle, and
composing them is an obvious thing to try. Route a token to a specialised expert,
then reuse \emph{that expert's} weights for as many steps as the token appears to
need. The resulting architecture inherits an efficiency argument from each side.
What it does not inherit is evidence: whether the composition beats either
component at a matched budget has not, to our knowledge, been isolated.

We isolate it. One training engine implements three configurations differing only
in the fields their definitions require: routing without recursion (MoE, six
experts, depth 1), recursion without routing (MoR, one shared block, depth up to
7), and the composition (MoRE, six experts, depth up to 7). All three sit at
$\approx$5.58M parameters, share a bit-identical input representation, and train
on the same fixed split of WikiText-103 \citep{merity2017wikitext} with the same
optimiser budget. Because the comparison is matched rather than merely similar, a
difference between arms cannot be a difference in the harness, the data, or the
parameter count.

\paragraph{What we find.} The composition does not pay off in the way its
premise predicts, and the reason is measurable rather than mysterious.

\begin{itemize}[leftmargin=1.4em,itemsep=1.5pt,topsep=2pt]
\item \textbf{Recursion alone predicts best in domain.} MoR reaches the lowest
  validation loss (3.515 nats/token) against MoRE's 3.637 and MoE's 3.948. Both
  recursive arms beat MoE; the composition lands between them.

\item \textbf{Recursion is genuinely adaptive; the composition is not.} MoR's
  learned halting policy beats \emph{every} uniform depth we can force, including
  its own best, by 0.54 nats. MoRE's learned policy merely ties the maximum
  forced depth, because its halting collapsed to always-run-to-budget: 97.3\% of
  validation tokens are forced out at the depth ceiling. This is the cleanest
  single contrast in the study, and it inverts the premise that adding routing
  would help the model target computation more precisely.

\item \textbf{Routing buys robustness, and efficiency, not in-domain quality.}
  Out of domain the ordering inverts: MoRE attains the best out-of-domain loss on
  the Pile \citep{gao2020pile}, ahead of MoR and well ahead of MoE. It does so at
  2.19$\times$ MoE's inference FLOPs against MoR's 9.07$\times$, which makes its
  gain over MoE roughly $3\times$ cheaper per nat than MoR's.

\item \textbf{The two mechanisms compete for one budget.} Depth and
  specialisation are not independent knobs in MoRE the way they are in MoR and
  MoE separately: a calibration sweep that recovers adaptive depth by raising the
  ponder penalty destroys expert differentiation, and the full-corpus arm that
  preserves differentiation saturates depth instead.

\item \textbf{One plausible mechanism is refuted.} The natural account---that
  routing confines a token to one expert's subspace, trapping the recursion---is
  falsified three ways: participation ratio is \emph{higher} in the recursive
  arms, per-step state velocity decays smoothly rather than vibrating, and
  successive state updates are near-orthogonal rather than rescaling.
\end{itemize}

We report this as a negative result about the composition at this scale, with the
mechanism attached, rather than as a claim about either mechanism in general. The
paper's contribution is the decomposition: because MoE \emph{is} MoRE with the
depth budget set to one and all three arms are parameter-matched from one
codebase, each contrast varies one thing. The scale at which the dilution account
predicts the composition would begin to pay is stated as a falsifiable claim in
the Limitations section.

\section{Related Work}
\label{sec:related}

\paragraph{Sparse expert routing.}
Conditional mixtures of local experts date to \citet{jacobs1991moe}; the
sparsely-gated form that made them practical at scale is due to
\citet{shazeer2017moe}, with Top-1 dispatch, load-balancing auxiliaries and the
capacity-factor machinery developed in GShard \citep{lepikhin2021gshard}, Switch
Transformer \citep{fedus2022switch} and ST-MoE \citep{zoph2022stmoe}. The standard
justification is more parameters at roughly constant per-token FLOPs. That
argument is unavailable here by construction: we pin all three arms to the same
parameter count, so what remains of the sparsity case is a quality and robustness
case, which \S\ref{sec:results} tests directly.

\paragraph{Adaptive depth and weight-tied recursion.}
Adaptive Computation Time \citep{graves2016act} introduced the learned per-step
halting unit we use; Universal Transformers \citep{dehghani2019ut} applied it to
a weight-tied transformer block, and PonderNet \citep{banino2021pondernet}
reformulated the halting distribution to keep a gradient on the halt head near
saturation. Early-exit variants place the decision at the layer or token level
instead \citep{elbayad2020depthadaptive,schuster2022calm,raposo2024mod}. Weight
tying itself is well established as a parameter-reduction device
\citep{lan2020albert}. What this literature reports is overwhelmingly
\emph{mean} step counts and exit-rate histograms. \S\ref{sec:halting} argues that
neither statistic tests input-dependence, and that the comparison which does---against
the best uniform depth at the same budget---is rarely made.

\paragraph{Composing the two.}
Recent work has begun to combine routing with recursion, most directly by using a
router to assign per-token recursion depth \citep{bae2025mor,raposo2024mod}.
Those systems route \emph{to depths} over shared weights. MoRE differs in that
routing selects among genuinely independent expert FFNs and the recursion then
reuses \emph{that expert's} weights, so the two decisions are made by different
heads over different objects and a token keeps one expert for its whole
trajectory. We are not aware of a matched-budget comparison of such a composition
against each of its components, which is the gap this paper fills. We make no
novelty claim for the combination itself.

\paragraph{Evaluation at small scale.}
Minimal-pair benchmarks \citep{warstadt2020blimp} and likelihood-based world
knowledge probes \citep{leong2023ewok} score a \emph{relative} preference within a
matched pair, so they degrade gracefully with model scale rather than collapsing
to chance the way multiple-choice knowledge suites do. At 5.6M parameters they are
the appropriate instrument; \S\ref{sec:results} reports them with the caveat that
our vocabulary is fitted to the training corpus and absolute accuracies are
therefore not comparable to published models.
